\chapter{Closure of loop ileostomy}
 
\section*{Definition and Overview}
Closure of a loop ileostomy, often termed ileostomy reversal, is a surgical procedure to restore small bowel continuity after a temporary loop ileostomy. A loop ileostomy is typically created during an abdominal surgery (such as a total mesorectal excision for rectal cancer, or subtotal colectomy for ulcerative colitis) to protect a distal bowel join (anastomosis) while it heals. The reversal surgery involves dissecting the loop of the ileum from the surrounding abdominal wall layers, closing the opening in the bowel wall (using sutures or staples, or via segmental resection and anastomosis), and returning the small intestine back into the peritoneal cavity before closing the abdominal wall fascia and skin.
 
\section*{Epidemiology}
Loop ileostomy is the most common temporary defunctioning stoma used in colorectal surgery. Reversal is typically scheduled 8 to 12 weeks after the initial stoma creation, depending on the patient's nutritional status, recovery, and any postoperative adjuvant therapies. Overall morbidity after reversal ranges from 15\% to 40\%, with small bowel obstruction and surgical site infections being the most frequent complications. While it is generally a shorter and less invasive procedure than the primary surgery, it carries a small but important risk of readmission.
 
\section*{Aetiology and Pathophysiology}
Reversal of a loop ileostomy is indicated once the clinical need for small bowel diversion has resolved:
 
\begin{itemize}
    \item \textbf{Healed distal join:} confirmation that the downstream colorectal, ileorectal, or ileoanal anastomosis is fully intact, healed, and patent.
    \item \textbf{Absence of stenosis or leaks:} ensuring no stricture or leakage exists in the downstream segment that could cause obstruction or sepsis.
    \item \textbf{Surgical procedure:} the surgeon makes an incision around the stoma, dissects it free from abdominal wall muscles, and performs either a primary closure (suturing or stapling the enterotomy) or a resection of the stoma-bearing segment with end-to-end anastomosis.
\end{itemize}
 
\section*{Clinical Features}
Postoperative clinical signs reflect recovery and the return of normal bowel movements:
 
\begin{itemize}
    \item Resumption of normal gastrointestinal motility, indicated by the passage of gas (flatus) and stool within 2 to 4 days
    \item Mild to moderate abdominal cramps and soreness at the surgical site
    \item Alterations in bowel frequency, urgency, and loose consistency (stools may be frequent and watery in the first few weeks as the bowel adapts)
    \item Perianal skin irritation from frequent or loose bowel movements
    \item General postoperative fatigue and recovery of physical strength
\end{itemize}
 
\section*{Differential Diagnosis}
Postoperative issues must be distinguished from major complications:
 
\begin{itemize}
    \item \textbf{Postoperative ileus:} a temporary delay in bowel function, characterised by abdominal bloating, nausea, and vomiting without mechanical blockage.
    \item \textbf{Adhesive small bowel obstruction (SBO):} mechanical obstruction of the small bowel due to early adhesions or kinking of the anastomosis.
    \item \textbf{Anatomotic leak:} a critical complication resulting in pelvic or abdominal sepsis, presenting with severe pain, fever, peritonitis, and tachycardia.
    \item \textbf{High-output stoma syndrome:} (if reversal fails or stoma is recreated) causing profound dehydration.
\end{itemize}
 
\section*{When to Seek Medical Care}
Patients must monitor their recovery and report any warning signs.
 
\textbf{Call 000 (or local emergency services) for:}
\begin{itemize}
    \item Sudden, severe abdominal pain or progressive swelling of the abdomen
    \item Persistent vomiting or inability to keep fluids down
    \item Signs of sepsis: high fever, shaking, rapid heart rate, or rapid breathing
    \item Dizziness, confusion, or fainting (suggesting severe dehydration or internal bleeding)
\end{itemize}
 
\section*{Diagnosis and Investigations}
\begin{itemize}
    \item \textbf{Preoperative studies:} essential to perform a water-soluble contrast enema (gastrografin study) or flexible endoscopy to verify the healing and patency of the downstream bowel before surgery.
    \item \textbf{Clinical evaluation:} monitoring vitals, fluid intake and output, and wound status.
    \item \textbf{Blood tests:} full blood count, kidney function, and serum electrolytes (critical to monitor for dehydration and electrolyte depletion).
    \item \textbf{Abdominal imaging:} plain X-ray or CT scan if bowel obstruction, ileus, or anastomotic leak is suspected postoperative.
\end{itemize}
 
\section*{Management}
\begin{itemize}
    \item \textbf{Surgical technique:} mobilisation and closure of the ileum, frequently using a purse-string skin closure to reduce wound infection rates.
    \item \textbf{Enhanced Recovery After Surgery (ERAS):} protocols including early oral intake, minimal intestinal spasms, and early mobilisation to promote gut motility.
    \item \textbf{Fluid and electrolyte management:} close monitoring of hydration, as fluid loss can occur rapidly due to changes in intestinal transit.
    \item \textbf{Bowel management:} use of bulking agents (e.g., psyllium husk) or antidiarrhoeals (e.g., loperamide) once downstream patency is confirmed, to slow transit and reduce stool frequency.
\end{itemize}
 
\section*{Complications}
\begin{itemize}
    \item Small bowel obstruction (SBO) or postoperative ileus
    \item Anastomotic leak or enterocutaneous fistula
    \item Surgical site infection (SSI) at the peristomal wound
    \item Dehydration and acute kidney injury from high stool volume
    \item Incisional hernia at the reversal site
\end{itemize}
 
\section*{Prognosis and Outcomes}
The prognosis after closure of a loop ileostomy is generally excellent, with most patients restoring normal bowel function and returning to work and normal activities within 4 to 6 weeks. While temporary changes in bowel habits (increased frequency and urgency) are very common initially, these symptoms usually improve significantly over 3 to 12 months as the ileum and colon adapt. A successful reversal restores intestinal continuity permanently in the vast majority of patients.
 
\section*{Prevention and Self-Care}
\begin{itemize}
    \item Keep the surgical wound clean and dry, following specific dressing instructions, and check daily for signs of infection
    \item Drink plenty of fluids (water, electrolyte solutions) to prevent dehydration, especially if stools are loose or frequent
    \item Eat a low-fibre (low-residue) diet for the first 2 to 4 weeks to avoid overloading the healing small bowel, and introduce high-fibre foods slowly
    \item Protect perianal skin by using barrier creams if experiencing frequent bowel movements
    \item Avoid heavy lifting (over 5 kg) and strenuous physical activity for 6 weeks to reduce the risk of a hernia at the wound site
\end{itemize}
 
\section*{Sources and Further Reading}
The following reputable organisations provide further information for healthcare students and patients seeking reliable, current advice about ileostomy reversal.
 
\begin{itemize}
    \item Association of Coloproctology of Great Britain and Ireland. \textit{Patient guide to ileostomy reversal.} https://www.acgbbi.org.uk/
    \item Healthdirect Australia. \textit{Ileostomy reversal: what to expect during recovery.} https://www.healthdirect.gov.au/
    \item Mayo Clinic. \textit{Ileostomy: reversal surgery, risks, and recovery.} https://www.mayoclinic.org/
    \item NHS. \textit{Ileostomy: stoma reversal process.} https://www.nhs.uk/conditions/ileostomy/reversal/
    \item MSD Manual. \textit{Jejunostomy and ileostomy: management and reversal.} https://www.msdmanuals.com/
\end{itemize}
